\section{Mass as latency: the Fibonacci resolution coordinate (interface)}
\label{sec:mass_latency_coordinate}

This section records the resolution-depth coordinate and the operational mass-as-delay dictionary used at the matching layer.
These definitions are used downstream by the closed mass-spectrum template and by the falsifiability statements formulated in the protocol language.

\paragraph{A readout viewpoint (interface).}
In a static-ontology reading (Remark~\ref{rem:static_ontology}), the underlying protocol substrate is fixed while physical episodes reflect sequential access and finite-resolution projection.
On this viewpoint, increasing effective resolution does not change the substrate; it changes which structure is rendered/accessible in the readout, and how much overhead (latency) is paid to stabilize it.

\subsection{Fibonacci resolution coordinate}
\label{subsec:mass_resolution_coordinate}

Fix the reference scale $\mu_0=m_e$ and define the resolution coordinate
\begin{equation}\label{eq:r_of_mu_z128}
r(\mu):=\frac{\log(\mu/m_e)}{\log\varphi},
\end{equation}
where $\varphi=(1+\sqrt{5})/2$ is the golden ratio.
Equivalently, the exponential map is
\begin{equation}\label{eq:mu_of_r_z128}
\mu(r)=m_e\,\varphi^{\,r}.
\end{equation}
This is the resolution-flow dictionary used throughout the paper.

\paragraph{Tick-first meaning.}
In the tick-first dictionary (Section~\ref{sec:tick_calculus}), mass and energy are not primitive substances: they are calibrated names for \emph{time-scale ratios}.
Accordingly, the coordinate $r(\mu)$ is used as a single log-time coordinate: by the Compton-clock relation, it is simultaneously a log-frequency coordinate and (up to sign) a log-time coordinate (Remark~\ref{rem:mass_as_compton_clock} and Appendix~\ref{app:time_mass_delay}).
This is the precise sense in which the present framework treats ``mass as depth'' as ``mass as a time-lag / overhead coordinate'' rather than as an independent input.

\subsubsection{Mass as delay: scattering time lag as inertia (interface)}
\label{subsubsec:mass_as_delay_interface}
In the scan-based identification dictionary used in this paper, ``mass as depth'' is read as ``mass as protocol overhead'': deeper stabilization requires additional local protocol resources.
An operational proxy for such overhead is measurable scattering delay.
Let $S(\omega)$ be a (nearly) unitary scattering matrix at angular frequency $\omega$.
The Wigner--Smith time-delay matrix is
\[
Q(\omega):=-\iu\,S(\omega)^\dagger\frac{\dd S}{\dd \omega},
\qquad
\tau_{\mathrm{WS}}(\omega):=\Tr Q(\omega),
\]
and in a one-channel setting $S(\omega)=\e^{\iu\delta(\omega)}$ one has $\tau_{\mathrm{WS}}(\omega)=\dd\delta/\dd\omega$ (Section~\ref{app:time_mass_delay}).

Why should depth create mass?
Physically, high information density forces the one-dimensional scan to execute additional local folding and consistency operations to stabilize a persistent pattern.
When probed through a scattering channel, this manifests as an excess time delay.
To an external observer, a localized excitation that systematically ``lags'' in response behaves as an inertial degree of freedom.
Throughout this paper this dictionary is used only at the matching layer: it complements the Compton-clock relation below and provides an independent falsifiability route via delay-derived lapse ratios (Section~\ref{subsec:p6_wigner_smith_delay} and Section~\ref{app:time_mass_delay}).

\paragraph{Micro/macro unification pointer.}
\AuditTag For the unified wormhole/black-hole contract form reused across scales (micro trapping as dwell-time proxy; paid pointer-jump shortcuts under hard gates; uplift as the bridge carrier),
see Section~\ref{subsec:micro_bh_wh_particle_unification}, Section~\ref{subsec:x3a_wormhole_rigidity_unified_constraints}, and
Proposition~\ref{prop:four_way_unification_contract_form}.

\begin{remark}[Mass as a clock rate (matching dictionary)]\label{rem:mass_as_compton_clock}
By the standard relations $E=mc^2$ and $E=\hbar\omega$ \cite{Einstein1905,NielsenChuang2010}, a mass scale $\mu$ defines a Compton angular frequency $\omega_C(\mu)=\mu c^2/\hbar$ and a Compton time scale $\tau_C(\mu)=1/\omega_C(\mu)=\hbar/(\mu c^2)$.
Therefore the resolution coordinate \eqref{eq:r_of_mu_z128} is equivalently a log-frequency (or log-time) coordinate:
\[
r(\mu)=\log_\varphi\!\left(\frac{\omega_C(\mu)}{\omega_C(m_e)}\right)
=-\log_\varphi\!\left(\frac{\tau_C(\mu)}{\tau_C(m_e)}\right).
\]
In particular, the depth mismatch $\Delta r=r-\widehat r$ reported later in the mass-spectrum closure is the same multiplicative mismatch in Compton-clock period, $\tau_C/\tau_{C,\mathrm{pred}}=\varphi^{-\Delta r}$, and can be compared to operational delay proxies (Section~\ref{app:time_mass_delay}) at the matching layer.
The same logarithmic coordinate also linearizes Schwarzschild black-hole thermodynamics: $\log_\varphi S_{\mathrm{BH}}=2\,r(M)+\text{const}$ and $\log_\varphi T_H=-\,r(M)+\text{const}$ (Appendix~\ref{app:bh_wormholes_pointer}).
\end{remark}

\begin{remark}[Why the base $\varphi$ is canonical on the golden branch]\label{rem:phi_base_entropy}
Lemma~\ref{lem:xm_fib} gives $|X_m|=F_{m+2}$ for the admissible language at window length $m$.
By Binet's formula, $F_{m+2}$ grows exponentially as $\varphi^{m}$ up to a fixed prefactor \cite{HardyWright2008}.
Equivalently, the golden-mean shift has topological entropy $\log\varphi$ and its Artin--Mazur zeta function is $\zeta(z)=1/(1-z-z^2)$ (Lemma~\ref{lem:golden_zeta_matrix}) \cite{LindMarcus1995}.
Thus using $\log\varphi$ as the denominator in \eqref{eq:r_of_mu_z128} matches the intrinsic exponential growth rate of the admissible symbolic language on the golden branch.
\end{remark}


